% Part: first-order-logic
% Chapter: natural-deduction
% Section: proof-theoretic-notions

\documentclass[../../../include/open-logic-section]{subfiles}

\begin{document}

\iftag{FOL}
      {\olfileid{fol}{ntd}{ptn}}
      {\olfileid{pl}{ntd}{ptn}}

\olsection{ثبوتی-نظریاتی تصورات (Proof-Theoretic Notions)}

\begin{editorial}
ایہ حصہ فطری استنباط لئی ثابت پذیری دے تعلق تے بے تضادی دیاں
تعریفاں اکٹھیاں کردا اے۔
\end{editorial}

\begin{explain}
جیویں اسیں کئی اہم معنوی تصورات (منطقی اعتبار، منطقی نتیجہ تے پورا ہون
جوگتا) دی تعریف دِتّی اے، ہُن اسیں اوہناں دے مطابق
\emph{ثبوتی-نظریاتی تصورات (proof-theoretic notions)} دی تعریف دیندے آں۔
ایہناں دی تعریف !!{sentence}s دے !!{structure}s وچ پورا اترن ول رجوع کر کے
نہیں، بلکہ کجھ !!{sentence}s دا ہور جملیاں توں !!{derivability}
یا !!{nonderivability} ول رجوع کر کے دِتّی جاندی اے۔ ایہ اک اہم دریافت سی
کہ ایہ تصور آپو وچ میل کھاندے نیں۔ اوہناں دے میل کھان دی گل
\emph{صحت مندی (soundness)} تے \emph{تمامیت دے قضیاں (completeness theorems)}
دا مضمون اے۔
\end{explain}


\begin{defn}[قضیے (Theorems)]
!!{sentence}~$!A$ اک \emph{قضیہ (theorem)} اے جے فطری استنباط وچ
~$!A$ دی کوئی !!a{derivation} ہووے جیہدے سارے مفروضے !!{discharged} ہون۔
اسیں $\Proves !A$ لکھدے آں جے $!A$ قضیہ اے، تے جے ایہ قضیہ نہیں تاں
$\Proves/ !A$ لکھدے آں۔
\end{defn}

\begin{defn}[!!^{derivability}]
!!^a{sentence} $!A$، !!{sentence}s دے سیٹ~$\Gamma$ توں
\emph{!!{derivable}} اے، $\Gamma \Proves !A$، جے کوئی
!!{derivation} ایہو جہی ہووے جیہدا نتیجہ~$!A$ ہووے تے جیہدا ہر مفروضہ
یا تاں !!{discharged} ہووے یا~$\Gamma$ وچ ہووے۔ جے $!A$، $\Gamma$ توں
!!{derivable} نہیں تاں اسیں $\Gamma \Proves/ !A$ لکھدے آں۔
\end{defn}

\begin{defn}[بے تضادی (Consistency)]
!!{sentence}s دا سیٹ~$\Gamma$ \emph{متضاد (inconsistent)} اے اُدوں تے صرف
اُدوں جدوں $\Gamma
\Proves \lfalse$۔ جے $\Gamma$ متضاد نہیں، یعنی جے
$\Gamma \Proves/ \lfalse$، تاں اسیں اوہنوں \emph{بے تضاد (consistent)}
آکھدے آں۔
\end{defn}

\begin{prop}[انعکاسیت (Reflexivity)]
\ollabel{prop:reflexivity}
جے $!A \in \Gamma$ اے، تاں $\Gamma \Proves !A$۔
\end{prop}

\begin{proof}
مفروضہ $!A$ اپنے آپ وچ~$!A$ دی اک !!a{derivation} اے جیہدے سارے
!!{undischarged} مفروضے (یعنی $!A$)~$\Gamma$ وچ نیں۔
\end{proof}
  

\begin{prop}[یک رُخیت (Monotonicity)]
\ollabel{prop:monotonicity}
جے $\Gamma \subseteq \Delta$ تے $\Gamma \Proves !A$، تاں $\Delta
\Proves !A$۔
\end{prop}

\begin{proof}
~$!A$ دی~$\Gamma$ توں ہر !!{derivation}،~$!A$ دی~$\Delta$ توں وی
!!a{derivation} اے۔
\end{proof}

\begin{prop}[تعدیت (Transitivity)]
\ollabel{prop:transitivity}
جے $\Gamma \Proves !A$ تے $\{!A\} \cup \Delta \Proves
!B$، تاں $\Gamma \cup \Delta \Proves !B$۔
\end{prop}

\begin{proof}
جے $\Gamma \Proves !A$ اے، تاں اک !!a{derivation}~$\delta_0$،~$!A$
دی اے، جیہدے سارے !!{undischarged} مفروضے~$\Gamma$ وچ نیں۔ جے $\{!A\} \cup
\Delta \Proves !B$ اے، تاں اک !!a{derivation}~$\delta_1$،~$!B$
دی اے، جیہدے سارے !!{undischarged} مفروضے~$\{!A\} \cup \Delta$ وچ نیں۔
ہُن ایہ ویکھو:
\begin{prooftree}
  \AxiomC{$\Delta, \Discharge{!A}{1}$}
  \RightLabel{$\delta_1$}
  \DeduceC{$!B$}
  \DischargeRule{\Intro{\lif}}{1}
  \UnaryInfC{$!A \lif !B$}
  \AxiomC{$\Gamma$}
  \RightLabel{$\delta_0$}
  \DeduceC{$!A$}
  \RightLabel{\Elim{\lif}}
  \BinaryInfC{$!B$}
\end{prooftree}
ہُن سارے !!{undischarged} مفروضے $\Gamma \cup
\Delta$ وچ نیں؛ ایس لئی ایہ وکھاؤندا اے کہ $\Gamma \cup \Delta \Proves !B$۔
\end{proof}

جدوں $\Gamma = \{!A_1, !A_2, \ldots, !A_k\}$ کوئی محدود سیٹ ہووے، تاں اسیں مختصر علامت $!A_1,!A_2,\ldots,!A_k \Proves !B$، $\Gamma \Proves !B$ دی تھاں ورت سکدے آں؛ خاص طور تے $!A \Proves !B$ دا مطلب اے کہ $\{!A\} \Proves !B$۔

دھیان دیو کہ جے $\Gamma \Proves !A$ تے $!A
\Proves !B$، تاں $\Gamma \Proves !B$۔ ایہ وی نکلدا اے کہ جے $!A_1,
\dots, !A_n \Proves !B$ تے $\Gamma \Proves !A_i$ ہر~$i$ لئی، تاں
$\Gamma \Proves !B$۔

\begin{prop}
\ollabel{prop:incons}
ہِٹھلیاں گلاں اک دوجے دے برابر نیں:
    \begin{enumerate}
      \item \( \Gamma \) متضاد اے۔
      \item \( \Gamma \Proves {!A} \)، ہر !!{sentence}~\( {!A} \) لئی۔
      \item \( \Gamma \Proves {!A} \) تے \( \Gamma \Proves \lnot {!A} \)، کسے !!{sentence}~\( {!A} \) لئی۔
    \end{enumerate}
\end{prop}

\begin{proof}
مشق۔
\end{proof}

\tagprob{FOL}
\begin{prob}
\olref[fol][ntd][ptn]{prop:incons} ثابت کرو۔
\end{prob}
\tagendprob

\tagprob{notFOL}
\begin{prob}
\olref[pl][ntd][ptn]{prop:incons} ثابت کرو۔
\end{prob}
\tagendprob

\begin{prop}[محدود گواہی (Compactness)]
  \ollabel{prop:proves-compact}
  \begin{enumerate}
  \item جے $\Gamma \Proves !A$، تاں کوئی محدود ذیلی سیٹ $\Gamma_0
    \subseteq \Gamma$ ایہو جہا اے کہ $\Gamma_0 \Proves !A$۔
  \item جے~$\Gamma$ دا ہر محدود ذیلی سیٹ
    بے تضاد اے، تاں $\Gamma$ بے تضاد اے۔
  \end{enumerate}
\end{prop}

\begin{proof}
  \begin{enumerate}
    \item جے $\Gamma \Proves !A$، تاں
      اک !!a{derivation}~$\delta$،~$!A$ دی~$\Gamma$ توں اے۔ منّ لو $\Gamma_0$،~$\delta$ دے
      !!{undischarged} مفروضیاں دا سیٹ اے۔ چونکہ ہر !!{derivation} محدود
      ہوندی اے، ایس لئی~$\Gamma_0$ وچ صرف محدود گنتی دے
      !!{sentence}s ہو سکدے نیں۔ سو~$\delta$،~$!A$ دی اک !!a{derivation} محدود~$\Gamma_0 \subseteq \Gamma$
      توں اے۔
    \item ایہ (1) دا متقابلِ عکس اے، خاص حالت $!A
      \ident \lfalse$ لئی۔
  \end{enumerate}
\end{proof}

\end{document}
